% !TEX root = ../main.tex
\section{Results and Discussion (Preliminary)}
\label{sec:results}

\irtwo{Results and Discussion (preliminary). For Interim Report \#2 this section needs: (1) the architecture as it stands (context diagram, aircraft IBD, the three requirements tagged \texttt{improved} as the measured gap) --- the figures are now in Section~\ref{sec:system-model}; (2) the two hand checks from \texttt{experiment-definition.md} \S9 written up as preliminary results; (3) the first OpenAP numbers --- Figure~\ref{fig:fuelflow-altitude} and Table~\ref{tab:outcomes} are in place, the discussion is not; (4) one paragraph on what is not yet done.}

\subsection{Model Results}
\plantodo{S10, S13 (D-028)}{\textbf{Model results.} With the architecture now presented in Section~\ref{sec:system-model}, report its measures here, against the measures named in Section~\ref{sec:deliverables}: for example, how many stakeholder requirements trace to a kept objective, how many system requirements are allocated to a constituent system, and which requirements are tagged \texttt{improved} because today's NAS does not meet them. Without this, the systems engineering work is described but never assessed.}

\subsection{Experiment Results}
\label{sec:experiment-results}

% The table body is generated by simulation/experiment_scoring.py (do not edit the
% file it inputs). It replaces the hand-typed table whose cells were empty.
\begin{table}[h]
\centering
\small
\caption{Outcome of each case at four levels, for the base scenario scored at the actual masses. Each cost is a difference from flying on undisturbed. Preliminary: the scenario's numbers are not all confirmed.}
\label{tab:outcomes}
% Generated by simulation/experiment_scoring.py; do not edit. Fidelity: actualMass
% Scenario: base; run 2026-10-29
\begin{tabular}{@{}llp{3.0cm}p{3.2cm}p{2.6cm}p{2.4cm}@{}}
\toprule
Case & Option & Largest cost to any one aircraft (kg fuel-equiv.) & Airline totals (USD) & Sector (calls / effort / risk) & System fuel, CO2 (kg) \\
\midrule
A & O10 & 40 (AC2) & X 0, Y 102, Z 0 & 1+1 / 1 / low & 41, 130 \\
B & O8 & 26 (AC1) & X 66, Y 0, Z 0 & 1+1 / 3 / low & 32, 102 \\
C & O4 & 12 (AC2) & X 0, Y 32, Z 0 & 1+1 / 3 / low & 12, 39 \\
D & O4 & 12 (AC2) & X 0, Y 32, Z 0 & 1+1 / 3 / low & 12, 39 \\
\bottomrule
\end{tabular}

\end{table}

\plantodo{S12}{\textbf{Which option each case chose, and why.} Walk Table~\ref{tab:outcomes} row by row. For case A, name the rule that produced the choice. For B, C and D, name the cost that decided it.}

\plantodo{S12}{\textbf{Did the climb rule matter?} With the crossing pair above FL350, cases A and B could not offer a climb and cases C and D could. Say whether case C used one, for which aircraft, and what it saved. If C did not climb either, say so: that is a result too.}

\plantodo{S12}{\textbf{B minus A, and C minus B.} The first is what automation is worth; the second is what the information exchange is worth. Compare the size of each with the prior figure of tens of gallons per flight (\cite{engility2014tasarAlaska}, p.~12) and say whether the difference is larger than the truth model's uncertainty.}

\plantodo{S9, S12}{\textbf{Does the best option change with the level?} Read across each row of Table~\ref{tab:outcomes}. Name any option that is better for one aircraft or one airline and worse for the sector or the system, and who pays (the largest cost to any one aircraft). This is the sub-system versus system discussion.}

\plantodo{S12}{\textbf{Airline X's two flights.} Discuss the option that touches both of one airline's aircraft: what the airline would choose between its late flight and its early one, and that the controller cannot see that preference. Connect to the dispatcher priorities (schedule integrity and connections rated at or above fuel; source of that table still to confirm).}

\plantodo{S12}{\textbf{Workload is three things.} Discuss clearances issued, effort to devise, and residual risk separately. Say where a ``more workload'' option on one of them is ``less'' on another.}

\plantodo{S12}{\textbf{When would the controller say no?} (owner's note, second hand review, 5.2). Discuss whether refusals are more likely in morning rushes at hub airports or near holidays, using the sector-loading rule in \cite{engility2014tasarAlaska} (p.~10). A hypothesis; needs a source.}

\plantodo{S12}{\textbf{Is schedule the right priority?} (owner's note, second hand review, 4.11). Contrast a carbon price, which makes fuel weigh more against time, with controllers' usual practice of giving schedule the priority.}

\begin{figure}[h]
\centering
\figorbox{b739_fuelflow_altitude.png}{PLACEHOLDER: run \texttt{simulation/b739\_envelope.py} to produce \texttt{report/figures/b739\_fuelflow\_altitude.png}}
\caption{Fuel flow and specific range against altitude at Mach 0.80 for the heavy and light B737-9 (OpenAP).}
\label{fig:fuelflow-altitude}
\end{figure}

\plantodo{Track A}{\textbf{Figure slot: sequence diagram} of case C's conflict resolution (controller, ATC decision support, airline operations center, FMS, crew), rendered from the model once drafted (turn 3).}

\plantodo{S12}{\textbf{The case where nothing changes.} Report the variation with similar weights and the climb closed, where case C picks the same option as case B, to show the comparison is not rigged.}

Each column of Table~\ref{tab:outcomes} belongs to a part of the context diagram (Figure~\ref{fig:sysml-systemContext}). The aircraft column belongs to the aircraft and its flight crew. The airline column belongs to the airline operations center within flight operations. The sector column belongs to air traffic control within airspace management. The system column belongs to no single part. It is a property of the NAS as a whole, which is the reason no constituent system computes it today and the reason the model is needed to state it.

The exchange that cases C and D add is one interface in the model: the airline operations center's performance share to air traffic control's decision support, carrying gross weight, cost index and the flight's delay margin. It satisfies system requirement SLR-INF-02, which is tagged as part of the improved architecture because today's NAS does not meet it. The interface is not drawn on Figure~\ref{fig:sysml-systemContext}, which shows the baseline.


\subsection{Emissions and Who Pays}
\label{sec:emissions-impact}

\plantodo{S9, S12}{\textbf{Airline level: put CO2 into dollars.} Take the fuel difference between cases from Table~\ref{tab:outcomes}, convert to kg of CO2, then to dollars. State the route used and its source. Candidates: (1) the price of a carbon credit or allowance, if a scheme applies to the flight; (2) passenger value, meaning revenue gained or lost because some passengers choose the flight listed with lower emissions. Make the point that fuel is already counted once as fuel cost and that pricing its CO2 counts the same kilograms again, on purpose.}

\plantodo{S9, S12}{\textbf{Passenger value.} Discuss how passengers see a flight's emissions when booking, and whether that reaches the airline as demand. The author's own practice of choosing lower-emission listings is the motivation, not the evidence: cite a study. Note the gap between one decision in one sector and the label a passenger saw at booking; the link is the airline's average over many flights.}

\plantodo{S12}{\textbf{Scale it honestly.} One conflict in one sector is a small number. If a per-year or fleet-wide figure is offered, state the multiplier and its source (for example, how often level crossing conflicts occur), and label it illustrative.}

\plantodo{S9}{\textbf{Who pays.} For the option the system prefers, name the aircraft or airline that is worse off and by how much. This is the equity discussion.}
